%% kapa: exam=1 round=32 year=2021 period=1 form=A subject=real_estate subject_ko=부동산학원론 seq=27 number=107
%% source: Q-Net 감정평가사 기출문제 (cst00309) artlSeq=5208397 file=extracted/1008/5208397/A형/2021년도 제32회 감정평가사 1차 1교시 A형.pdf
%% conversion: kordoc 4.24.51 -> tools/gen1.py (auto); review=proofread-sonnet; flags=none
\begin{question}{107}
토지개발방식으로서 수용방식과 환지방식의 비교에 관한 설명으로 옳지 않은 것은? (단, 사업구역은 동일함)
\begin{choices}
  \item 수용방식은 환지방식에 비해 종전 토지소유자에게 개발이익이 귀속될 가능성이 큰 편이다.
  \item 수용방식은 환지방식에 비해 사업비의 부담이 큰 편이다.
  \item 수용방식은 환지방식에 비해 기반시설의 확보가 용이한 편이다.
  \item 환지방식은 수용방식에 비해 사업시행자의 개발토지 매각부담이 적은 편이다.
  \item 환지방식은 수용방식에 비해 종전 토지소유자의 재정착이 쉬운 편이다.
\end{choices}
\end{question}
\answer{1}
